\section{Explanation and falsifiable predictions}
\label{sec:hypotheses}
The candidate explanation concerns a fixed training recipe: increasing the width of a trainable feature basis can accelerate learning of shared arithmetic structure while leaving a substantial interval after training-set fit. It predicts an earlier generalization event at larger width, rather than merely earlier memorization. This is an explanation to test, not an established capacity law. Changing width also changes initialization and optimization dynamics; the experiment does not isolate parameter count from those changes.

The prospective protocol was fixed after exploratory regime-finding runs and before the confirmatory seed blocks. Its three predictions have distinct evidential requirements. Let $T_m$ and $T_g$ denote confirmed training-fit and held-out generalization onsets, and let $R_g=\min(T_g,H)$ at the common horizon $H$; Section~\ref{sec:events} specifies their operational definitions.

\paragraph{P1: earlier restricted generalization.}
Under the main recipe, the paired mean contrast
\begin{equation}
\Delta_g=\mathbb{E}[R_g(512)-R_g(32)]
\label{eq:primary}
\end{equation}
is negative. A confidence interval wholly below zero supports the directional prediction; an interval wholly above zero contradicts it. An interval spanning zero, or pervasive joint censoring, leaves it unresolved. The full five-width response is reported even if it is nonmonotone.

\paragraph{P2: delay persists at the larger width.}
More than half of width-512 runs exhibit confirmed generalization after fit with both $T_g-T_m\geq1000$ updates and $T_g/T_m\geq10$. These cutoffs define a substantial delay for this study, rather than a natural boundary of grokking. A 95\% Wilson lower bound above $0.5$ supports P2; an upper bound below $0.5$ contradicts it; otherwise it remains unresolved. Continuous delays, censored lower bounds and threshold sensitivities accompany this classification.

\paragraph{P3: feature freezing attenuates the width advantage.}
For the same first eight seed blocks, define $D_c=\mathbb{E}[R_g(512,c)-R_g(32,c)]$, where $c$ is full training or frozen hidden features. The prediction is $D_{\mathrm{full}}-D_{\mathrm{frozen}}<0$. The representation-versus-fitting interpretation additionally requires at least six of eight frozen runs at \emph{each} endpoint width to reach training fit; the corresponding full-model counts must also be reported. Failure of this gate leaves that interpretation unresolved, even if the restricted-clock interaction has the predicted sign. Freezing changes both hidden adaptation and hidden weight shrinkage, so this intervention cannot separate their contributions.

Alternative outcomes matter: width may accelerate only fitting, the response may depend on the optimizer recipe, or frozen features may exhibit a comparable width advantage. Support for the complete account requires P1, P2 and an interpretable P3. Mixed outcomes warrant a conditional or revised explanation while retaining the original predictions. Neither a width sweep alone nor a functional Fourier diagnostic identifies a unique internal circuit.
